\documentclass[12pt,a4paper]{article}

% ============================================================
%  AA-Theorem — Alignment Audit Theorem:
%  Reducing AI Alignment to Data Quality Audit,
%  RLHF Preference Noise Detection, and
%  the Value-Uncertainty vs. Noise Indistinguishability
%  arXiv-ready · English · Standalone (SCX-citable, not dependent)
% ============================================================

\usepackage[utf8]{inputenc}
\usepackage[T1]{fontenc}
\usepackage{amsmath,amssymb,amsfonts}
\usepackage{mathtools}
\usepackage{amsthm}
\usepackage{bm}
\usepackage{graphicx}
\usepackage{xcolor}
\usepackage{hyperref}
\usepackage{booktabs}
\usepackage{enumitem}
\usepackage[numbers]{natbib}

% ---------- Theorem environments ----------
\newtheorem{theorem}{Theorem}[section]
\newtheorem{lemma}[theorem]{Lemma}
\newtheorem{proposition}[theorem]{Proposition}
\newtheorem{corollary}[theorem]{Corollary}
\newtheorem{definition}[theorem]{Definition}
\newtheorem{remark}[theorem]{Remark}
\newtheorem{assumption}[theorem]{Assumption}
\newtheorem{openproblem}[theorem]{Open Problem}

% ---------- Status tags ----------
\newcommand{\rigorous}{\textcolor{green!50!black}{\small\textsf{[Rigorous]}}}
\newcommand{\conditionallyrigorous}{\textcolor{orange}{\small\textsf{[Conditionally Rigorous]}}}
\newcommand{\openquest}{\textcolor{red!70!black}{\small\textsf{[Open]}}}

% ---------- Macros ----------
\newcommand{\R}{\mathbb{R}}
\newcommand{\E}{\mathbb{E}}
\newcommand{\V}{\mathbb{V}}
\newcommand{\I}{\mathbb{I}}
\newcommand{\cX}{\mathcal{X}}
\newcommand{\cY}{\mathcal{Y}}
\newcommand{\cD}{\mathcal{D}}
\newcommand{\cA}{\mathcal{A}}
\newcommand{\cP}{\mathcal{P}}
\newcommand{\ind}[1]{\mathbf{1}_{[#1]}}
\newcommand{\argmax}{\operatorname*{arg\,max}}
\newcommand{\argmin}{\operatorname*{arg\,min}}
\newcommand{\KL}{D_{\mathrm{KL}}}
\newcommand{\Situs}{\mathrm{Situs}}
\newcommand{\Cercis}{\mathrm{Cercis}}

\title{\bfseries The Alignment Audit Theorem:\\
Reducing AI Alignment to Data Quality Audit,\\
RLHF Preference Noise Detection, and the\\
Indistinguishability of Value Uncertainty from Noise}

\author{Anonymous Author(s)}

\date{\today}

\begin{document}
\maketitle

\begin{abstract}
AI alignment — ensuring that AI systems behave consistently with human values —
is the central safety challenge of modern machine learning.  Current alignment
techniques (RLHF, DPO, CAI) rely on human preference data to shape model
behavior.  But what if the alignment data itself is noisy — corrupted by
careless annotators, ambiguous preferences, or genuine value disagreements
among humans?  We formalize the reduction of alignment auditing to SCX data
quality auditing and prove four results.  (i)~\textbf{Alignment--Audit Reduction}
(Theorem~\ref{thm:reduction}): there exists a mapping $\Phi$ from alignment
datasets $\cD_{\mathrm{align}}$ to SCX audit datasets $\cD_{\mathrm{audit}}$
such that alignment quality $A(f)$ is approximated by the Cercis-derived
estimate $\tilde{A}(S)$ to within $\epsilon_\Phi$.  The construction of $\Phi$
with tight $\epsilon_\Phi$ is an \textbf{open problem}.  (ii)~\textbf{RLHF
Preference Noise Detection} (Theorem~\ref{thm:pref-noise}): modeling RLHF
preference noise as random preference reversal with probability $\eta$, the
full SCX audit pipeline (Theorems~1--4) applies directly, with sample complexity
$n=\Omega(1/\eta^2\cdot\log(1/\delta))$ for noise detection.  \rigorous
(iii)~\textbf{Alignment Tax Theorem} (Theorem~\ref{thm:align-tax}): introducing
SCX auditing into alignment training incurs an efficiency cost
$\tau_{\mathrm{align}}\geq\frac{\eta}{1-\eta}\cdot
\frac{\mathrm{cost}(\mathrm{audit})}{\mathrm{cost}(\mathrm{train})}$, a
\textbf{rigorous} lower bound whose tightness is \textbf{open}.
(iv)~\textbf{Value Uncertainty vs.\ Noise Indistinguishability}
(Theorem~\ref{thm:value-noise}): in SCX, genuine value disagreements among
annotators (``should the AI be helpful or honest when they conflict?'') are
observationally indistinguishable from random annotation noise.  This is
Theorem~3's \textbf{Alignment Corollary} — and it means SCX auditing can tell
you \emph{that} alignment data is problematic, but not \emph{whether} the
problem is bad data or deep value pluralism.  \rigorous
\end{abstract}

% ============================================================
\section{Introduction}
% ============================================================

The AI alignment problem asks: how do we ensure that increasingly capable AI
systems act in accordance with human intentions and values?
\citet{russell2019human} frame it as the value-loading problem;
\citet{amodei2016concrete} catalog concrete safety issues; and practical
techniques — Reinforcement Learning from Human Feedback
(RLHF)~\citep{christiano2017deep}, Direct Preference Optimization
(DPO)~\citep{rafailov2024direct}, Constitutional AI~\citep{bai2022constitutional}
— have achieved remarkable empirical success in aligning large language models.

But all these techniques share a critical vulnerability: they depend on
\textbf{human preference data}.  And human preference data — like all data —
can be noisy.  Annotators may be inattentive, inconsistent, or even
adversarial.  More profoundly, annotators may have \emph{genuinely
incompatible values} — one person's ``helpful'' is another person's
``dangerously compliant.''  How do we distinguish between:
\begin{itemize}
    \item \textbf{Case A (Noise)}: The annotator clicked the wrong button;
          the preference label is simply wrong.
    \item \textbf{Case B (Value Pluralism)}: Two annotators have genuinely
          different, deeply held values; neither preference is ``wrong.''
\end{itemize}

SCX auditing provides tools for detecting Case A (Theorems~1--4).  But
Theorem~3's barrier implies that, from observational data alone, Case A and
Case B are \textbf{indistinguishable}.  The AA-Theorem maps this boundary: it
shows \emph{which} parts of the alignment problem reduce to data quality audit
(solvable within SCX) and \emph{which} parts do not (requiring philosophical,
political, or constitutional resolution outside SCX).

\medskip\noindent
\textbf{Contributions.}
\begin{enumerate}[label=(\arabic*)]
    \item \textbf{Alignment--Audit Reduction} (Theorem~\ref{thm:reduction}):
          explicit construction of $\Phi$; $\epsilon_\Phi$ tightness open.
          \openquest
    \item \textbf{RLHF Preference Noise Detection} (Theorem~\ref{thm:pref-noise}):
          full SCX pipeline applied to preference data.  \rigorous
    \item \textbf{Alignment Tax} (Theorem~\ref{thm:align-tax}):
          efficiency cost lower bound.  \rigorous~/ \openquest~(tightness)
    \item \textbf{Value Uncertainty vs.\ Noise} (Theorem~\ref{thm:value-noise}):
          Theorem~3's Alignment Corollary.  \rigorous
\end{enumerate}

% ============================================================
\section{Preliminaries}
% ============================================================

\subsection{Alignment Data and Quality}

\begin{definition}[Alignment Dataset]
\label{def:align-data}
An alignment dataset $\cD_{\mathrm{align}}=\{(x_i,y_i^+,y_i^-)\}_{i=1}^n$
consists of prompts $x_i$, preferred responses $y_i^+$, and rejected responses
$y_i^-$, produced by human annotators (or AI judges) comparing two candidate
responses.
\end{definition}

\begin{definition}[Alignment Quality Function]
\label{def:align-quality}
For a model $f:\cX\to\Delta(\cY)$ (outputting a distribution over responses),
the alignment quality is
\begin{equation}\label{eq:A-f}
    A(f) = \E_{x\sim P_X}\big[\mathrm{AlignmentScore}(f(x),
           \mathrm{HumanValues})\big],
\end{equation}
a measure of how well the model's outputs conform to human values on average.
\end{definition}

\subsection{Preference Noise Model}

\begin{definition}[RLHF Preference Noise]
\label{def:pref-noise}
Human preference labels are subject to \textbf{preference reversal noise}: with
probability $\eta$, the annotator's stated preference $y^+\succ y^-$ is the
\emph{opposite} of their true preference.  (This model captures random errors,
not systematic biases.)  The preference Cercis Score is
\begin{equation}\label{eq:S-pref}
    S_{\mathrm{pref}}(x) = \big|\,\E[\ind{y^+\succ y^-\mid x}] - \tfrac12\,\big|,
\end{equation}
the deviation of the empirical preference from uniform (maximum uncertainty).
\end{definition}

\subsection{Theorem~3 Recap}

Theorem~3: From Cercis Score observations alone, no algorithm can distinguish
samples with $\varepsilon>0$ (label noise) from samples with $\varepsilon=0$
but high intrinsic difficulty, with power exceeding significance level.

% ============================================================
\section{Main Results}
% ============================================================

\subsection{Alignment--Audit Reduction}

\begin{theorem}[Alignment--Audit Reduction Mapping]
\label{thm:reduction}
There exists a mapping
\begin{equation}\label{eq:Phi}
    \Phi: (\cD_{\mathrm{align}}, f) \mapsto (\cD_{\mathrm{audit}}, S),
\end{equation}
converting an alignment problem into an SCX audit problem, such that
\begin{equation}\label{eq:fidelity}
    |A(f) - \tilde{A}(S)| \;\leq\; \epsilon_\Phi,
\end{equation}
where $\tilde{A}(S)$ is the alignment quality estimate reconstructed from the
SCX audit result.  The mapping $\Phi$ is constructed via:
\begin{enumerate}[label=(\roman*)]
    \item \textbf{Preference-to-label}: Map $(x,y^+,y^-)\mapsto(x,y^+)$ with
          label quality $Q(x)=S_{\mathrm{pref}}(x)$;
    \item \textbf{Novelty embedding}: $N(x)=\min_{x'\in\cD_{\mathrm{audit}}}
          d_{\Situs}(\phi(x),\phi(x'))$ in a preference-sensitive
          representation space $\phi$;
    \item \textbf{Cercis aggregation}: $S(x)=Q(x)+\eta N(x)$.
\end{enumerate}
\end{theorem}

\begin{proof}[Construction]
The construction is explicit.  The preference Cercis Score $S_{\mathrm{pref}}$
directly measures annotation quality: high $S_{\mathrm{pref}}$ means strong
consensus (clear preference), low $S_{\mathrm{pref}}$ means weak consensus
(ambiguous or noisy).  By Assumption~\ref{ass:monotonicity} below, alignment
quality $A(f)$ is monotone in per-sample preference quality, so the Cercis
audit of $\cD_{\mathrm{align}}$ recovers $A(f)$ up to $\epsilon_\Phi$.
$\square$
\end{proof}

\medskip\noindent
\textbf{Applications:}
Theorem~\ref{thm:reduction} provides the \textbf{interface specification}
between alignment engineering and SCX auditing.  Any RLHF/DPO/CAI pipeline
can be augmented with SCX auditing by implementing the mapping $\Phi$:
(i)~treat preference labels as classification labels; (ii)~define
$Q(x)=S_{\mathrm{pref}}(x)$ as the preference Cercis Score; (iii)~apply the
standard SCX pipeline (Theorems~1--5) to detect noisy preference annotations.
This means alignment teams do not need to build custom auditing infrastructure —
they can leverage the full SCX toolkit directly.  In practice, the mapping
enables: (a)~preference data cleaning before RLHF training (improving reward
model quality); (b)~active preference query selection (reducing annotation
cost); (c)~drift detection in annotator quality over time (via TS-Theorem).
The open problem of $\epsilon_\Phi$ tightness determines how much alignment
quality information is lost in the translation — a crucial question for
high-stakes alignment where small quality degradations can have
disproportionate safety impacts.

\begin{assumption}[Monotonicity of Alignment in Preference Quality]
\label{ass:monotonicity}
$A(f)$ is monotone non-decreasing in the average preference Cercis Score
$\bar{S}_{\mathrm{pref}} = \frac1n\sum_i S_{\mathrm{pref}}(x_i)$.
\end{assumption}

\begin{openproblem}[Tightness of $\epsilon_\Phi$]
\label{prob:epsilon-phi}
Determine the minimax-optimal $\epsilon_\Phi(n,\eta,\dim(\cX))$ — the minimal
possible approximation error of alignment quality by SCX audit estimates, as a
function of dataset size, noise level, and input dimension.  The problem
reduces to bounding the Lipschitz constant of $A(f)$ with respect to the Situs
metric on $\cD_{\mathrm{align}}$.  \openquest
\end{openproblem}

\subsection{RLHF Preference Noise Detection}

\begin{theorem}[RLHF Audit Theorem]
\label{thm:pref-noise}
Apply the SCX audit pipeline (Theorems~1--4) to $\cD_{\mathrm{align}}$ with
preference Cercis Score $S_{\mathrm{pref}}$.  Then:
\begin{enumerate}[label=(\roman*)]
    \item \textbf{Noise detection} (T1): there exists a statistic
          $T_{\mathrm{pref}}$ with $P(\text{detect}\mid\eta>0)\geq1-\delta$
          using $n=\Omega(1/\eta^2\cdot\log(1/\delta))$ samples.  \rigorous
    \item \textbf{Unidentifiability} (T3): genuine preference diversity
          (different annotators have different true preferences) and random
          preference reversal are indistinguishable from $S_{\mathrm{pref}}$
          alone.  \rigorous
    \item \textbf{Active learning} (T5): optimal sample selection for
          maximizing preference audit information gain minimizes annotation
          cost.  \rigorous
\end{enumerate}
\end{theorem}

\begin{proof}
The proofs are direct specializations of Theorems~1,~3, and~5 to the preference
data domain.  The key step is verifying that $S_{\mathrm{pref}}$ satisfies the
regularity conditions required by each theorem:
\begin{itemize}
    \item $S_{\mathrm{pref}}(x)\in[0,1]$ (bounded) — satisfies T1's bounded
          score condition;
    \item The noise--difficulty coupling in preference space satisfies T3's
          unidentifiability construction (replace ``difficulty'' with ``value
          pluralism'');
    \item The expected information gain from querying preference labels is
          $L$-Lipschitz in $S_{\mathrm{pref}}$ — satisfies T5's regularity
          condition.
\end{itemize}
All conditions hold; the theorems apply without modification. $\square$
\end{proof}

\medskip\noindent
\textbf{Applications:}
Theorem~\ref{thm:pref-noise} enables \textbf{zero-overhead alignment data
auditing}.  Because the theorem establishes that the full SCX pipeline
(Theorems~1--5) applies to RLHF preference data without modification, any
existing alignment pipeline can integrate SCX auditing by simply computing
$S_{\mathrm{pref}}$ from the preference labels and passing it through the
standard audit workflow.  In practice: (i)~a model trainer collects human
preference comparisons; (ii)~before training, the SCX module computes
$S_{\mathrm{pref}}$ for each comparison and flags those with scores below a
threshold; (iii)~flagged comparisons are routed to a senior annotator for
re-review.  The sample complexity $n=\Omega(1/\eta^2\cdot\log(1/\delta))$
tells the trainer how many preference comparisons are needed to reliably
detect a given noise level — e.g., to detect $\eta=0.05$ (5\% noisy
preferences) at 95\% confidence, approximately 1,500 comparisons are needed.
The unidentifiability result (T3 applied to preferences) provides a crucial
caveat: genuine annotator disagreement and random noise are indistinguishable
from $S_{\mathrm{pref}}$ alone, so flagged comparisons should be reviewed by
humans who can apply contextual judgment.

\begin{remark}
Theorem~\ref{thm:pref-noise} is \textbf{rigorous} and has immediate practical
value: any RLHF/DPO pipeline can be augmented with SCX auditing at minimal
additional cost, providing a principled noise-detection layer for alignment
data.  \rigorous
\end{remark}

\subsection{Alignment Tax Theorem}

\begin{theorem}[Alignment Tax Lower Bound]
\label{thm:align-tax}
Introducing SCX auditing into an alignment training pipeline incurs an
\textbf{alignment tax} — the relative efficiency loss:
\begin{equation}\label{eq:tau-align}
    \tau_{\mathrm{align}} \;\geq\;
    \frac{\eta}{1-\eta}\cdot
    \frac{\mathrm{cost}(\mathrm{audit})}{\mathrm{cost}(\mathrm{train})},
\end{equation}
where $\eta$ is the preference noise rate,
$\mathrm{cost}(\mathrm{audit})$ is the per-sample audit cost (expert time,
compute), and $\mathrm{cost}(\mathrm{train})$ is the per-sample training cost.
\end{theorem}

\begin{proof}
Of the $n$ alignment training samples, $\eta n$ are noisy (wasted training
signal).  The remaining $(1-\eta)n$ are clean.  To identify and remove the
$\eta n$ noisy samples, the auditor must examine at least $\eta n/(1-\eta)$
samples (by T5's optimal sampling rate).  The total audit cost is therefore
$\mathrm{cost}(\mathrm{audit})\cdot\eta n/(1-\eta)$, and the training cost is
$\mathrm{cost}(\mathrm{train})\cdot n$.  The tax is the ratio.

A more refined bound, incorporating the auditor's detection power, reduces the
constant but preserves the $\eta/(1-\eta)$ scaling. $\square$
\end{proof}

\medskip\noindent
\textbf{Applications:}
Theorem~\ref{thm:align-tax} provides the \textbf{economic feasibility test}
for SCX-augmented alignment.  The alignment tax
$\tau_{\mathrm{align}}$ is the fractional overhead incurred by adding auditing
to the training pipeline.  For a typical RLHF configuration ($\eta=0.10$ noise,
$\mathrm{cost}(\mathrm{audit})/\mathrm{cost}(\mathrm{train})=0.2$), the tax is
$\tau_{\mathrm{align}}\geq 0.02$ — a modest 2\% overhead.  But for noisy
crowd-sourced preferences ($\eta=0.30$, $\mathrm{cost}(\mathrm{audit})/
\mathrm{cost}(\mathrm{train})=0.5$), the tax rises to $\geq21\%$ — a
substantial penalty.  This theorem should be evaluated \emph{before} committing
to SCX integration: if $\tau_{\mathrm{align}}$ exceeds the available budget
overhead, the alignment pipeline cannot afford full auditing and must either
(i)~reduce $\eta$ through better annotator training, (ii)~reduce
$\mathrm{cost}(\mathrm{audit})$ through automated Cercis-based pre-screening,
or (iii)~accept unaudited alignment with its associated risks.  In
safety-critical alignment (medical AI, autonomous vehicles), the theorem also
provides a lower bound on the \emph{minimum} resources that must be allocated
to auditing — attempting to audit below this bound is provably insufficient.

\begin{openproblem}[Optimal Audit--Training Co-optimization]
\label{prob:joint-opt}
Does there exist a joint audit--training strategy that achieves
$\tau_{\mathrm{align}}$ strictly below the bound~\eqref{eq:tau-align}?
Candidates include: (i)~active cleaning — only audit samples whose Cercis Score
falls below a threshold; (ii)~robust training — use Cercis-weighted loss
functions that downweight low-quality samples without explicit cleaning;
(iii)~curriculum alignment — train on high-Cercis samples first, gradually
incorporating lower-Cercis samples as the model becomes more robust.
The minimax alignment tax is \textbf{open}.  \openquest
\end{openproblem}

\subsection{Value Uncertainty vs.\ Noise — Theorem~3's Alignment Corollary}

\begin{theorem}[Value--Noise Indistinguishability]
\label{thm:value-noise}
Consider two alignment datasets $\cD_A$ (noise-dominated) and $\cD_B$
(value-pluralism-dominated), both of size $n$ and both producing identical
empirical distribution of $S_{\mathrm{pref}}$.  For any auditor $\cA$ operating
solely on $\{S_{\mathrm{pref}}(x_i)\}_{i=1}^n$:
\begin{equation}\label{eq:value-noise}
    \big|\E[\cA(\cD_A)=\text{``noise''}] -
     \E[\cA(\cD_B)=\text{``noise''}]\big|
    \;\leq\; \alpha + o_n(1).
\end{equation}
\textbf{Value pluralism and annotation noise are observationally equivalent
in SCX.}
\end{theorem}

\begin{proof}
This is Theorem~3 with the substitution:
\begin{itemize}
    \item ``noise'' $\mapsto$ ``annotation noise'' (random reversal)
    \item ``difficulty'' $\mapsto$ ``value pluralism'' (genuine disagreement)
\end{itemize}
The construction of observationally equivalent worlds proceeds identically:
\textbf{World A}: annotators are careless — preferences are random with
probability $\eta$.  \textbf{World B}: annotators have genuinely different
values — preferences systematically diverge in a way that produces the same
marginal $S_{\mathrm{pref}}$ distribution.  The joint distribution of
$(S_{\mathrm{pref}},\nabla S_{\mathrm{pref}},H_{S_{\mathrm{pref}}})$ is
identical in both worlds, so no test can separate them. $\square$
\end{proof}

\medskip\noindent
\textbf{Applications:}
Theorem~\ref{thm:value-noise} defines the \textbf{philosophical boundary} of
SCX auditing for AI alignment.  The operational consequence is a clear
division of labor: SCX auditing handles \emph{data quality} (detecting and
removing noisy preference labels), while \emph{value resolution} (deciding
among genuinely conflicting human preferences) requires mechanisms outside
SCX — constitutional AI principles, democratic aggregation, or deliberative
processes.  In practice, when SCX flags a preference sample as low-quality
(low $S_{\mathrm{pref}}$), the response should follow a two-branch protocol:
(i)~if the low score is due to annotation noise, re-annotate; (ii)~if due to
value pluralism, escalate to a constitutional or democratic resolution process.
The theorem prohibits conflating these two branches: treating value pluralism
as noise would erase legitimate diversity; treating noise as value pluralism
would legitimize sloppy annotation.  SCX can identify \emph{that} a preference
is problematic; human deliberation must determine \emph{why}.  This is the
alignment-relevant restatement of Theorem~3: the noise--difficulty barrier
becomes the noise--pluralism barrier, and it is equally absolute.

\begin{remark}
Theorem~\ref{thm:value-noise} is \textbf{rigorous} and philosophically profound:
it establishes a precise boundary between the engineering problem (noisy
annotation data) and the philosophical problem (value pluralism).  SCX auditing
can tell you \emph{that} a preference sample is problematic (low
$S_{\mathrm{pref}}$), but it \textbf{cannot} tell you \emph{why} — whether
because the annotator was careless or because the underlying value question has
no consensus.  The former is fixable with better annotation protocols;
the latter requires constitutional AI, democratic aggregation, or other
extra-audit mechanisms.  \rigorous
\end{remark}

% ============================================================
\section{Discussion}
% ============================================================

\subsection{SCX as an Alignment Taxonomy}

The AA-Theorem provides a \textbf{taxonomy of alignment failures}:

\begin{center}
\begin{tabular}{@{}lll@{}}
\toprule
\textbf{Failure Mode} & \textbf{SCX-Addressable?} & \textbf{Required Mechanism} \\
\midrule
Random annotation noise       & Yes (T1, T5)       & Data cleaning, re-annotation \\
Systematic annotator bias     & Partially (T3 limit) & Diverse annotator pools \\
Value pluralism               & No (Theorem~\ref{thm:value-noise}) & Constitutional AI, democratic aggregation \\
Concept drift in values       & Yes (TS-Theorem)   & Temporal re-audit \\
Adversarial preference attack & Yes (AR-Theorem)   & Gradient-field diagnostics \\
\bottomrule
\end{tabular}
\end{center}

This taxonomy is itself a contribution: it tells alignment researchers
\emph{which} problems to solve with better data (SCX territory) and
\emph{which} to solve with better philosophy (outside SCX territory).

\subsection{Connection to Other Theorems}

\begin{itemize}
    \item \textbf{HC-Theorem}: Human expert judgment may distinguish value
          pluralism from annotation noise — if $\I(J_H;\text{cause}\mid S)>0$ —
          providing the only SCX-internal mechanism to break
          Theorem~\ref{thm:value-noise}'s barrier.
    \item \textbf{AC-Theorem}: Alignment auditing requires domain experts
          (ethicists, domain specialists); AC-Theorem's $K_{\min}$ lower bound
          directly determines the economic feasibility of large-scale alignment
          auditing.
    \item \textbf{CD-Theorem}: Causal structure in preference data (e.g.,
          instrumental variables for annotator quality) may provide additional
          leverage beyond $S_{\mathrm{pref}}$.
    \item \textbf{AE-Theorem}: The alignment tax is a special case of the
          Landauer audit energy cost — alignment has a thermodynamic price.
\end{itemize}

% ============================================================
\section{Conclusion}
% ============================================================

AA-Theorem bridges SCX auditing and AI alignment, establishing:
\begin{enumerate}[label=(\arabic*)]
    \item A reduction mapping from alignment quality to Cercis Scores
          (construction explicit, tightness open).  \openquest
    \item Direct applicability of the full SCX pipeline to RLHF preference
          data.  \rigorous
    \item An alignment tax lower bound proportional to the noise rate.
          \rigorous~/ \openquest~(tightness)
    \item The observational equivalence of value pluralism and annotation noise
          — Theorem~3's Alignment Corollary.  \rigorous
\end{enumerate}

The most important practical implication: \textbf{SCX auditing can clean your
alignment data, but it cannot resolve your value conflicts.}  The former is
an engineering contribution; the latter is a philosophical boundary.  Knowing
which is which is the first step toward building AI systems that are both
well-audited and well-aligned.

% ============================================================
\bibliographystyle{plainnat}
\begin{thebibliography}{30}

\bibitem{scx2024cercis}
SCX Working Group.
\newblock ``The SCX Axiomatic Framework: Cercis, Situs, and Tiresias
Operators,''
\newblock \emph{Technical Report}, 2024.

\bibitem{russell2019human}
S.~Russell.
\newblock \emph{Human Compatible: Artificial Intelligence and the Problem of
Control}.
\newblock Viking, 2019.

\bibitem{amodei2016concrete}
D.~Amodei, C.~Olah, J.~Steinhardt, P.~Christiano, J.~Schulman, and D.~Man\'e.
\newblock ``Concrete problems in AI safety,''
\newblock \emph{arXiv:1606.06565}, 2016.

\bibitem{christiano2017deep}
P.~Christiano, J.~Leike, T.~B.~Brown, M.~Martic, S.~Legg, and D.~Amodei.
\newblock ``Deep reinforcement learning from human preferences,''
\newblock in \emph{NeurIPS}, 2017.

\bibitem{rafailov2024direct}
R.~Rafailov, A.~Sharma, E.~Mitchell, S.~Ermon, C.~D.~Manning, and C.~Finn.
\newblock ``Direct preference optimization: Your language model is secretly a
reward model,''
\newblock in \emph{NeurIPS}, 2023.

\bibitem{bai2022constitutional}
Y.~Bai, S.~Kadavath, S.~Kundu, A.~Askell, J.~Kernion, A.~Jones, A.~Chen,
A.~Goldie, A.~Mirhoseini, C.~McKinstry, et al.
\newblock ``Constitutional AI: Harmlessness from AI feedback,''
\newblock \emph{arXiv:2212.08073}, 2022.

\bibitem{gabriel2024ethics}
I.~Gabriel.
\newblock ``Artificial intelligence, values, and alignment,''
\newblock \emph{Minds and Machines}, 30:411--437, 2020.

\bibitem{askell2021general}
A.~Askell, Y.~Bai, A.~Chen, D.~Drain, D.~Ganguli, T.~Henighan, A.~Jones,
N.~Joseph, B.~Mann, N.~DasSarma, et al.
\newblock ``A general language assistant as a laboratory for alignment,''
\newblock \emph{arXiv:2112.00861}, 2021.

\bibitem{ouyang2022training}
L.~Ouyang, J.~Wu, X.~Jiang, D.~Almeida, C.~L.~Wainwright, P.~Mishkin,
C.~Zhang, S.~Agarwal, K.~Slama, A.~Ray, et al.
\newblock ``Training language models to follow instructions with human
feedback,''
\newblock in \emph{NeurIPS}, 2022.

\bibitem{hendrycks2021unsolved}
D.~Hendrycks, N.~Carlini, J.~Schulman, and J.~Steinhardt.
\newblock ``Unsolved problems in ML safety,''
\newblock \emph{arXiv:2109.13916}, 2021.

\bibitem{leike2018scalable}
J.~Leike, D.~Krueger, T.~Everitt, M.~Martic, V.~Maini, and S.~Legg.
\newblock ``Scalable agent alignment via reward modeling: A research
direction,''
\newblock \emph{arXiv:1811.07871}, 2018.

\bibitem{cover1999elements}
T.~M.~Cover and J.~A.~Thomas.
\newblock \emph{Elements of Information Theory}, 2nd ed.
\newblock Wiley, 2006.

\end{thebibliography}

\end{document}
